\documentclass{amsproc}
\newtheorem{theorem}{Theorem}
\title[Folded stops]{A spacing rule for folded stops}
\author{Fold Desk}\address{Model Gauge Workshop, Fold Bench}
\email{fold@example.invalid}\thanks{The fold bench supplied the trial strips.}
\author{Slot Desk}\address{Model Gauge Workshop, Slot Bench}
\email{slot@example.invalid}\thanks{The slot bench cut the test rail.}
\subjclass[2020]{Primary 05B45}
\keywords{paper gauge, spacing rule}
\begin{document}
\begin{abstract}
Two card stops leave a measurable gap in a model rail. A fixed fold preserves
this gap when the stops move together.
\end{abstract}
\maketitle
\section{A trial gauge}
\subsection{The common fold}
Cut two stops on the same grid as the strip in Figure~\ref{fig:rail}.
\begin{figure}[htbp]
\centering\rule{30mm}{2mm}
\caption{Rail for the folded stops}\label{fig:rail}
\end{figure}
\begin{theorem}
Sliding both stops by one grid step preserves their separation.
\end{theorem}
\begin{proof}
Let the original marks be $a$ and $b$. After a step of length $t$ the gap is
\begin{equation}\label{eq:shift}
(b+t)-(a+t)=b-a .
\end{equation}
\end{proof}
\newpage
\section{Bench check}
Table~\ref{tab:shift} checks equation~\eqref{eq:shift} against the workshop card~\cite{card}.
\begin{table}[htbp]
\centering\caption{Positions on the grid}\label{tab:shift}
\begin{tabular}{lrr}Trial & First & Second\\Before & 3 & 8\\After & 4 & 9\end{tabular}
\end{table}
\begin{thebibliography}{9}
\bibitem{card} Model Gauge Workshop, \emph{Grid rail card}, bench series.
\end{thebibliography}
\end{document}
